\scanpage{89}
% Typescript read after a half turn.
Signalons pour finir ces généralités, que la formation des duals $D(G)$ commute à l'extension de la base, qui transforme donc groupes réflexifs en groupes réflexifs.

Nous nous intéressons par la suite au cas où $\underline C=(\mathrm{Sch})_{/S}$, catégorie des préschémas sur $S$, et $I=\underline G_{m\ S}$, le ``groupe multiplicatif sur $S$'', cf exposé I. Pour tout groupe \deleted{commutatif} ordinaire $M$, nous considérons le $S$-groupe $M_S$. On voit aussitôt que pour tout schéma en \unsure{groupes|groupe} $I$ sur $S$, on a un isomorphisme canonique (fonctoriel en $M,I$ et compatible avec l'extension de la base):
\[
 \operatorname{Hom}_{S\text{-gr}}(M_S,I)=\operatorname{Hom}_{\mathrm{gr}}(M,I\ (I/S))
\]
Appliquant ceci à un $S'$ variable sur $S$, on trouve un isomorphisme \inserted{fonctoriel}
\[
 D(M_S)(S')\qquad \operatorname{Hom}_{\mathrm{gr}}(M,\underline G_m(S'))
\]
On retrouve le foncteur déjà considéré dans l'exposé I, \inserted{aussi noté $D(M)$,} qui est représentable \deleted{\unsure{par} \unsure{comme}} pour $M$ commutatif comme
\[
 D(M)=D(M_S)=\operatorname{Spec}\ \underline O_S\ M
\]
où $\underline O_S\ M$ désigne l'Algèbre du groupe $M$ à coéfficients dans $\underline O_S$. (No\deleted{\unsure{On le}}tons \inserted{d'ailleurs dans le cas général} que $D(M)$ ne change pas si on remplace $M$ par $M$ rendu \deleted{\uline{Définition} 1.1. Un} abélien, de sorte \deleted{\unsure{qu'il n'y a pas d'inconv}} qu'on ne perd rien en supposant $M$ commutatif).

\uline{Définition} 1.1. Un schéma en groupes \inserted{$G$} sur $S$ est dit diagonalisable s'il est isomorphe à un schéma de la forme \deleted{\unsure{$D($}} $D(M)=D(M_S)=\underline{\operatorname{Hom}}_{S\text{-gr}}(M_S,\underline G_{\underline m}$ pour un groupe commutatif $M$ convenable. On dit que $G$ est localement diagonalisable, si tout point de $S$ admet un voisinage ouvert $U$, tel que $G\ U$ soit diagonalisable.

\deleted{\uline{Définition} 1.2. Un schéma en groupes $I$ sur $S$ est appelé un \uline{réseau constant} (resp. un réseau \uline{localement constant}) s'il est isomorphe à un schéma en}
% archive sheet number: 88
